%%\usepackage[top=2cm,bottom=2cm,left=1cm,right=1cm]{geometry}


\begin{titlepage}
     \begin{center}
	\includegraphics[width=0.09\textwidth]{UNAM}\Large Universidad Nacional Autónoma de México
        	\includegraphics[width=0.09\textwidth]{FI}\\[1cm]
        \Large Facultad de Ingeniería\\[1cm]
       % \Large División de Ciencias Básicas\\[1cm]
         \Large Laboratorio de Fundamentos de Control(6655)\\[1cm]
         %la clave antes era:4314
         \footnotesize Profesor: Salcedo Ubilla María Leonor Ing.\\[1cm]
        \footnotesize Semestre 2019-1\\[1cm]
        
       

        \Large Práctica No. 1\\[1cm]
        
           

\Large Introdcción MATLAB
        
         %Texto a la derecha
          \begin{flushright}
\footnotesize  Grupo 2\\[0.5cm]
\footnotesize Brigada: 4\\[0.5cm]
\footnotesize Rodrigo Adrián Martínez López\\[0.5cm]
\footnotesize Vivar Colina Pablo\\[0.5cm]
 \end{flushright}
    %Texto a la izquierda
          \begin{flushleft}
        \footnotesize Ciudad Universitaria Agosto de 2018.\\
          \end{flushleft}
         
          
        %\vfill
        %\today
   \end{center}
\end{titlepage}
 %agregar portada

\documentclass{article}
\usepackage[utf8]{inputenc}
\usepackage[spanish.mexico]{babel}

\title{Dispositivos y circuitos Electrónicos previo 10}
\author{Pablo Vivar Colina}
\date{Abril 2018}

\usepackage{natbib}
\usepackage{graphicx}

\usepackage{tikz}
\usepackage[american voltages, american currents,siunitx]{circuitikz}

\begin{document}


\maketitle

%%\usepackage[top=2cm,bottom=2cm,left=1cm,right=1cm]{geometry}


\begin{titlepage}
     \begin{center}
	\includegraphics[width=0.09\textwidth]{UNAM}\Large Universidad Nacional Autónoma de México
        	\includegraphics[width=0.09\textwidth]{FI}\\[1cm]
        \Large Facultad de Ingeniería\\[1cm]
       % \Large División de Ciencias Básicas\\[1cm]
         \Large Laboratorio de Fundamentos de Control(6655)\\[1cm]
         %la clave antes era:4314
         \footnotesize Profesor: Salcedo Ubilla María Leonor Ing.\\[1cm]
        \footnotesize Semestre 2019-1\\[1cm]
        
       

        \Large Práctica No. 1\\[1cm]
        
           

\Large Introdcción MATLAB
        
         %Texto a la derecha
          \begin{flushright}
\footnotesize  Grupo 2\\[0.5cm]
\footnotesize Brigada: 4\\[0.5cm]
\footnotesize Rodrigo Adrián Martínez López\\[0.5cm]
\footnotesize Vivar Colina Pablo\\[0.5cm]
 \end{flushright}
    %Texto a la izquierda
          \begin{flushleft}
        \footnotesize Ciudad Universitaria Agosto de 2018.\\
          \end{flushleft}
         
          
        %\vfill
        %\today
   \end{center}
\end{titlepage}
 %agregar portada

%\section{Previo 8}







\section{Diodo Zener}

 Los diodos zener, zener diodo o simplemente zener, son diodos que están diseñados para mantener un voltaje constante en su terminales, llamado Voltaje o Tensión Zener (Vz) cuando se polarizan inversamente, es decir cuando está el cátodo con una tensión positiva y el ánodo negativa. Un zener en conexión con polarización inversa siempre tiene la misma tensión en sus extremos (tensión zener).\citep{dZener}\\

\begin{figure}[ht!]
   \centering
\includegraphics[scale=0.5]{Imagenes/diodo-zener.jpg}
\caption{}
     \label{fig:zener}
 \end{figure}
 
 \subsection{Cómo Funciona un Diodo Zener}

 Cuando lo polarizamos inversamente y llegamos a Vz el diodo conduce y mantiene la tensión Vz constante aunque nosotros sigamos aumentando la tensión en el circuito. La corriente que pasa por el diodo zener en estas condiciones se llama corriente inversa (Iz).
 Se llama zona de ruptura por encima de Vz. Antes de llegar a Vz el diodo zener NO Conduce.\citep{dZener}\\

 Como ves es un regulador de voltaje o tensión. Fijate en la gráfica de funcionamiento del zener más abajo.\citep{dZener}\\

 Cuando está polarizado directamente el zener se comporta como un diodo normal.\citep{dZener}\\

 Pero OJO mientras la tensión inversa sea inferior a la tensión zener, el diodo no conduce, solo conseguiremos tener la tensión constante Vz, cuando esté conectado a una tensión igual a Vz o mayor. Aquí puedes ver una la curva característica de un zener:\citep{dZener}\\
 
 \begin{figure}[ht!]
     \centering
     \includegraphics[scale=0.5]{Imagenes/curva-diodo-zener.jpg}
     \caption{Caption}
     \label{fig:curvaZener}
 \end{figure}
 

\section{Regulación de Voltaje}

Un regulador de tensión o regulador de voltaje es un dispositivo electrónico diseñado para mantener un nivel de tensión constante.\citep{regTensWiki}\\

Los reguladores electrónicos de tensión se encuentran en dispositivos como las fuentes de alimentación de los computadores, donde estabilizan las tensiones de Corriente Continua usadas por el procesador y otros elementos. En los alternadores de los automóviles y en las plantas generadoras, los reguladores de tensión controlan la salida de la planta. En un sistema de distribución de energía eléctrica, los reguladores de tensión pueden instalarse en una subestación o junto con las líneas de distribución de forma que todos los consumidores reciban una tensión constante independientemente de qué tanta potencia exista en la línea.\citep{regTensWiki}\\

\begin{figure}[ht!]
    \centering
    \includegraphics[scale=1]{Imagenes/50px-Voltage_Regulator.png}
    \caption{Terminales 	Entrada, tierra/ajuste y salida}
    \label{fig:simb}
\end{figure}

\section{Regulador de voltaje con diodo Zener}

\subsection{Características de los reguladores de voltaje con diodo zener}

El diodo Zener se puede utilizar para regular una fuente de voltaje. Este semiconductor se fabrica en una amplia variedad de voltajes y potencias. Estos van desde menos de 2 voltios hasta varios cientos de voltios, y la potencia que pueden disipar va desde 0.25 watts hasta 50 watts o más.\citep{regZener}\\

La potencia que disipa un diodo zener es simplemente la multiplicación del voltaje para el que fue fabricado por la corriente que circula por él. $P_z = (V_z)(I_z)$. Esto significa que la máxima corriente que puede atravesar un diodo zener es: $I_z = \frac{P_z}{V_z}$. (en amperios). Donde:\citep{regZener}\\

\begin{itemize}
    \item  $I_z$ = Corriente que pasa por el diodo Zener
     \item  $P_z$ = Potencia del diodo zener (dato del fabricante)
     \item $V_z$ = Voltaje del diodo zener (dato del fabricante)
\end{itemize}

Ejemplo: La corriente máxima que un diodo zener de 10 Voltios y 50 Watts puede aguantar, será: $I_z = \frac{P_z}{V_z}= \frac{50}{10} = 5$ amperios.\citep{regZener}\\


\begin{figure}[ht!]
    \centering
    \begin{circuitikz}
    

        \draw  (0,0) to[R,l=$R_s$](3,0); 
        \draw   (3,-3)to[zzDo,v=Vz](3,0)
        (0,-3)--(6,-3)
        (3,0)--(6,0)
        
        ;
        
        \node[draw] at (-0.4,0) {$V_{in}$};
        \node[draw] at (6.5,0) {$V_{out}$};
        
       
    \end{circuitikz}
    \caption{Rectificador Onda Completa}
    \label{fig:rectificadorOndaCompleta}
\end{figure}

\subsection{Cálculo de resistor limitador de corriente $R_s$}

El cálculo del resistor Rs está determinado por la corriente que pedirá la carga (lo que vamos a conectar a esta fuente de voltaje). Ver esquema del regulador de voltaje con diodo zener, con el resistor Rs conectado entre $V_{in}$ y el cátodo del zener. Este resistor se puede calcular con la siguiente fórmula: $R_s = (\frac{V_{enmin} – V_z}{1.1} ) IL_{max}$, donde:\\

\begin{itemize}
    \item     $V_{enmin}$: es el valor mínimo del voltaje de entrada. (acordarse que es un voltaje no regulado y puede variar)
    \item $IL _{max}$: es el valor de la máxima corriente que pedirá la carga.
\end{itemize}

Una vez conocido $R_s$, se obtiene la potencia máxima del diodo zener, con ayuda de la siguiente fórmula: $PD = (\frac{V_{enmin}-V_z}{R_s-IL_{min}})V_z$ \citep{regZener}\\


\subsubsection{Ejemplo de un diseño de regulador de voltaje con zener}

Una fuente de voltaje de 15 voltios debe alimentar una carga con 9 Voltios, que consume una corriente que varía entre 200 y 350 mA. (miliamperios). Se escoge un diodo zener de 9.1 voltios (muy cercano a 9 voltios:\\

\begin{itemize}
    \item Cálculo de $R_s: Rs =\frac{(15-9.1)}{(1.1×0.35)}= 15 \Omega$
\item Cálculo de la potencia del diodo zener: PD = (\frac{(15 – 9.1)}{15})9.1 = 3.58 [W]
\end{itemize}

Como no hay un diodo zener de 3.58 Vatios, se escoge uno de 5 vatios que es el más cercano.\\

\begin{itemize}
    \item Potencia de $R_s$: Un cálculo adicional es la potencia del resistor $R_s$. Este se hace con la fórmula: $P = I2 x R$. Ver Potencia en una resistencia (ley de Joule)
\end{itemize}

Los datos actuales son: I (max) = 350 miliamperios = 0.35 amperios y Rs = 15 Ohmios. Aplicando la fórmula, PRs = 0.352 x 15 = 1.84 Watts. Esto significa que a la hora de comprar este resistor deberá ser de 2 Watts o más.\citep{regZener}\\[2cm]



\bibliographystyle{plain}
\bibliography{referenciasPrev9}
\end{document}